\documentclass[12pt,a4paper]{article}

% ── Packages ──────────────────────────────────────────────
\usepackage{fontspec}
\usepackage[french]{babel}
\usepackage[margin=2.5cm]{geometry}
\usepackage{xcolor}
\usepackage{enumitem}
\usepackage{titlesec}
\usepackage{fancyhdr}
\usepackage{booktabs}
\usepackage{hyperref}
\usepackage{tcolorbox}
\usepackage{microtype}

% ── Colors ────────────────────────────────────────────────
\definecolor{primary}{RGB}{26, 86, 219}
\definecolor{accent}{RGB}{5, 150, 105}
\definecolor{warning}{RGB}{217, 119, 6}
\definecolor{gray}{RGB}{107, 114, 128}
\definecolor{lightgray}{RGB}{243, 244, 246}

\hypersetup{colorlinks=true, linkcolor=primary, urlcolor=primary, citecolor=primary}

% ── Section styling ───────────────────────────────────────
\titleformat{\section}
  {\Large\bfseries\color{primary}}
  {\thesection.}{0.6em}{}
  [\color{primary}\titlerule]

% ── Custom boxes ──────────────────────────────────────────
\tcbuselibrary{skins, breakable}

\newtcolorbox{infobox}[1][]{
  enhanced, breakable,
  colback=primary!8, colframe=primary,
  fonttitle=\bfseries, title={Info\quad #1},
  left=6pt, right=6pt, top=4pt, bottom=4pt, boxrule=1pt
}

\newtcolorbox{warningbox}[1][]{
  enhanced, breakable,
  colback=warning!10, colframe=warning,
  fonttitle=\bfseries, title={Attention\quad #1},
  left=6pt, right=6pt, top=4pt, bottom=4pt, boxrule=1pt
}

\newtcolorbox{questionbox}{
  enhanced, breakable,
  colback=accent!8, colframe=accent,
  left=6pt, right=6pt, top=4pt, bottom=4pt, boxrule=1pt
}

% ── Header / Footer ───────────────────────────────────────
\pagestyle{fancy}
\fancyhf{}
\fancyhead[L]{\small\color{gray} TP -- Segmentation d'Images Satellites avec U-Net}
\fancyhead[R]{\small\color{gray} Systemes Repartis \& GPU}
\fancyfoot[C]{\small\color{gray} \thepage}
\renewcommand{\headrulewidth}{0.4pt}

\begin{document}

% ── Title page ────────────────────────────────────────────
\begin{titlepage}
  \centering
  \vspace*{2cm}

  {\color{primary}\rule{\linewidth}{2pt}}
  \vspace{0.5cm}

  {\Huge\bfseries Questions du TP}\\[0.4cm]
  {\LARGE\bfseries\color{primary} Segmentation d'Images Satellites\\avec U-Net}

  \vspace{0.5cm}
  {\color{primary}\rule{\linewidth}{2pt}}

  \vspace{1.5cm}

  {\large
    \begin{tabular}{rl}
      \textbf{Module}       & Systemes Repartis et Programmation Parallele \\[4pt]
      \textbf{Filiere}      & Master S2 -- Intelligence Artificielle \\[4pt]
      \textbf{Universite}   & Faculte des Sciences Ben M'sick \\[4pt]
      \textbf{Groupe}       & GT019 \\[4pt]
      \textbf{Presentateurs}& Anass Dabibe \quad Hassan El Hadi \quad Marouane Ait Tamanait \\
    \end{tabular}
  }

  \vspace{2cm}

  \begin{tcolorbox}[colback=primary!5, colframe=primary, width=0.85\linewidth,
    boxrule=1pt, left=10pt, right=10pt, top=8pt, bottom=8pt]
    \centering
    \textbf{Consigne} \\[6pt]
    Executez le notebook \texttt{TP\_UNet\_Segmentation.ipynb} sur Kaggle (GPU),
    realisez les 6 captures d'ecran demandees, puis repondez aux 5 questions
    de reflexion ci-dessous dans un rapport PDF.
  \end{tcolorbox}

  \vfill
  {\color{gray}\small Annee universitaire 2025--2026}
\end{titlepage}

% ══════════════════════════════════════════════════════════
\section{Captures d'ecran a realiser}

Au fil de l'execution du notebook, realisez les \textbf{6 captures d'ecran}
suivantes (reperees par l'icone appareil photo dans le notebook) :

\begin{center}
\begin{tabular}{cll}
  \toprule
  \textbf{\#} & \textbf{Section du notebook} & \textbf{Contenu a capturer} \\
  \midrule
  1 & 1 -- Dataset            & Grille 4 images / 4 masques \\
  2 & 3 -- Entrainement       & Premieres et dernieres epoques (Loss, IoU) \\
  3 & 4 -- Courbes            & Courbes de Loss et d'IoU \\
  4 & 5 -- Evaluation         & Tableau d'IoU par classe \\
  5 & 6 -- Visualisation      & Figure de predictions (4 x 3) \\
  6 & 7 -- Benchmark          & Resultat CPU vs GPU et acceleration \\
  \bottomrule
\end{tabular}
\end{center}

\begin{infobox}
  Chaque capture doit etre accompagnee d'une breve description (2--3 phrases)
  expliquant ce que l'on observe.
\end{infobox}

% ══════════════════════════════════════════════════════════
\section{Questions de reflexion}

Repondez de maniere argumentee a chacune des questions suivantes.

\begin{questionbox}
  \begin{enumerate}[leftmargin=*, label=\textbf{Q\arabic*.}, itemsep=12pt]

    \item \textbf{Architecture} \\
          Pourquoi U-Net utilise-t-il des \textit{skip connections} ?
          Que se passerait-il si on les supprimait ?

    \item \textbf{Metrique} \\
          Pourquoi utilise-t-on l'IoU (\textit{Intersection over Union})
          plutot que l'accuracy pour evaluer la segmentation ?

    \item \textbf{Parallelisme} \\
          Quelle est la difference entre \texttt{DataParallel} et
          \texttt{DistributedDataParallel} (DDP) dans PyTorch ?
          Lequel est plus efficace et pourquoi ?

    \item \textbf{Resultats} \\
          Quelle classe obtient le meilleur IoU dans vos resultats ?
          Proposez une explication basee sur les caracteristiques
          visuelles de cette classe dans les images.

    \item \textbf{GPU} \\
          D'apres votre benchmark, quel est le facteur d'acceleration
          GPU obtenu ? Comment ce resultat justifie-t-il l'utilisation
          des systemes repartis pour l'entrainement de modeles de
          deep learning ?

  \end{enumerate}
\end{questionbox}

% ══════════════════════════════════════════════════════════
\section{Livrable attendu}

\begin{warningbox}[A deposer sur Google Classroom]
  Un fichier \textbf{PDF unique} nomme \texttt{TP\_GT0XX\_NomPrenom.pdf}.
\end{warningbox}

\begin{enumerate}[leftmargin=*, label=\textbf{\arabic*.}, itemsep=4pt]
  \item Les \textbf{6 captures d'ecran} demandees, chacune accompagnee
        d'une breve description (2--3 phrases).
  \item Les \textbf{reponses} aux 5 questions de reflexion (Q1 a Q5).
  \item La figure \textbf{predictions} en grand format
        (image satellite $|$ verite terrain $|$ prediction U-Net).
\end{enumerate}

\vspace{0.5em}
\begin{center}
  \begin{tabular}{lc}
    \toprule
    \textbf{Element}                & \textbf{Points} \\
    \midrule
    Captures d'ecran (6 x 1 pt)     & 6  \\
    Questions Q1--Q5 (5 x 2 pts)    & 10 \\
    Qualite de la redaction         & 4  \\
    \midrule
    \textbf{Total}                  & \textbf{20} \\
    \bottomrule
  \end{tabular}
\end{center}

% ══════════════════════════════════════════════════════════
\section{Ressources}

\begin{itemize}[leftmargin=*, itemsep=4pt]
  \item \textbf{Paper U-Net} : \url{https://arxiv.org/abs/1505.04597}
  \item \textbf{Dataset Kaggle} :
        \url{https://www.kaggle.com/datasets/humansintheloop/semantic-segmentation-of-aerial-imagery}
  \item \textbf{PyTorch DataParallel} :
        \url{https://pytorch.org/docs/stable/generated/torch.nn.DataParallel.html}
  \item \textbf{PyTorch DDP tutorial} :
        \url{https://pytorch.org/tutorials/intermediate/ddp_tutorial.html}
\end{itemize}

\vfill
\begin{center}
  \color{gray}\small
  Bon courage ! --- Groupe GT019 --- Systemes Repartis \& GPU 2025--2026
\end{center}

\end{document}
